\documentclass[11pt]{article}
\usepackage[a4paper,margin=1in]{geometry}
\usepackage{amsmath,amssymb,amsthm}

\theoremstyle{plain}
\newtheorem{theorem}{Theorem}[section]
\newtheorem{lemma}[theorem]{Lemma}
\newtheorem{proposition}[theorem]{Proposition}
\newtheorem{corollary}[theorem]{Corollary}
\theoremstyle{definition}
\newtheorem{definition}[theorem]{Definition}
\theoremstyle{remark}
\newtheorem{remark}[theorem]{Remark}

% ------------------------------------------------------------------
% Notation (Shoenfield, with code(u) for the encoding).
% ------------------------------------------------------------------
\newcommand{\code}[1]{\mathit{code}(#1)}
\newcommand{\Con}{\mathit{Con}}
\newcommand{\Fun}{\mathit{Fun}}
\newcommand{\Len}{\mathit{Len}}
\newcommand{\Defines}{\mathit{Defines}}
\newcommand{\Prr}{\mathit{Pr}}
\newcommand{\Thy}{\mathit{TH}}
\newcommand{\Thyo}{\mathit{TH}_\omega}
\newcommand{\thmT}{\mathbf{th}}
\newcommand{\sub}{\mathbf{sub}}
\newcommand{\sat}{\models}
\newcommand{\Eval}{\mathbf{Eval}}

\title{Krivine's models-inside-models proof of G\"odel~II,\\
and its translation to Guard's BRA\\
\large{(reading of \emph{Th\'eorie axiomatique des ensembles}, pp.~126--127)}}
\author{}
\date{}

\begin{document}
\maketitle

\begin{abstract}
This note records Krivine's \emph{semantic} (Berry-style) proof of the
second incompleteness theorem as it is actually written in
\emph{Th\'eorie axiomatique des ensembles} (pp.~126--127), and gives a
precise dictionary translating it into Guard's basic recursive
arithmetic~(BRA).  The point of interest is that the proof's engine is
not an abstract ``uniformity across all models'' but a concrete
\emph{models-inside-models} construction: a model $M\sat\Thy$, under the
contradiction hypothesis $\Thy\vdash\Con(\Thy)$, contains an internal
object $N'$ that $M$ believes is a model of $\Thy$, and the value of the
Berry function $\nu$ is pinned down by the internal soundness
equivalence ``$N'\sat F$ holds in $M$ iff $F\in\Thy$.''  We argue that
this is the von~Neumann ``copy-the-antinomy'' route.  After two
corrections recorded below --- Krivine's definability is \emph{not}
decidable (a model gives \emph{bivalence}, not decidability), and BRA's
quantifier-freeness forces the $\Sigma_1$ content out of the object
language --- we reach a stable conclusion: \textbf{Krivine's argument
\emph{does} render syntactically in Guard's system}, as a model-free
Berry proof in which definability is $\thmT$-provability ($\Sigma_1$), the
Berry function $\nu$ is non-computable and lives at the formula level, and
the hypothesis $\Thy\vdash\Con(\Thy)$ is spent on the $\Pi_1$
non-definability fact.  No relativised model and no decidable definability
are needed.  This supersedes the vague ``Lindenbaum cut'' of the companion
note \texttt{berry-krivine-design.tex}.
\end{abstract}

% =================================================================
\section{Provenance, and what is new here}
% =================================================================

The one-page summary \texttt{krivine.pdf} states the shape of Krivine's
argument (Berry definability, a counting bound, ``uniformity across
models'') but omits its mechanism.  The photographs \texttt{IMG\_0196},
\texttt{IMG\_0197} (and the duplicate \texttt{Unknown.jpeg}) reproduce
the full proof from Krivine's textbook, p.~126 and its continuation.
The full text reveals what carries the uniformity, and that is the new
material this note records.

\paragraph{Legibility caveat.}  These are photographs of small print
with sub/superscript formulas.  The \emph{structure} below is reported
with confidence; the precise exponent towers and the exact reading of
the internal theoremhood predicate ``$T(F)^{\circ}\omega$'' / the
auxiliary theory $\Thyo$ are reported as they appear but flagged where a
symbol is not fully legible.  None of the BRA translation depends on
those exact constants.

% =================================================================
\section{Krivine's proof as written (pp.~126--127)}
% =================================================================

\subsection{The theorem}

\begin{theorem}[Krivine, p.~126]
Let $\Thy$ be an axiomatised theory containing $ZF$.  If $\Thy$ is
consistent, then ``$\neg\,\Con(\Thy)$'' is also consistent --- i.e.\
$\Con(\Thy)$ is not a consequence of $\Thy$.
\end{theorem}

The proof is by contradiction: \emph{assume every universe (model)
satisfying $\Thy$ also satisfies $\Con(\Thy)$}, i.e.\ assume
$\Thy\vdash\Con(\Thy)$.

\subsection{The Berry function $\nu$, internally}

Krivine introduces $p=\nu(x)$ \emph{by an internal formula with two
free variables} expressing:
\begin{quote}
``$p$ is the least integer that, \emph{in at least one model $M$
of~$\Thyo$}, is not definable by a parameter-free one-variable formula
$F(x)$ of length~$<x$.''
\end{quote}
Two sub-formulas make this internal, and they are the crux:
\begin{description}
\item[``$M$ is a model of $\Thyo$''] is written
  \[
     \forall F\;\bigl(F\text{ closed}\ \wedge\ T(F)^{\circ}\omega
        \ \rightarrow\ M\sat F\bigr),
  \]
  i.e.\ for every closed formula provable in $\Thyo$, $M$ satisfies it.
  (Here $T(\cdot)^{\circ}\omega$ is Krivine's internal theoremhood
  predicate; the exact decoration is not fully legible, but its role is
  ``$F$ is a theorem.'')
\item[``$\nu(x)$ is defined in $M$ by $F$''] is the internal
  definability predicate $\Defines$: $F$ is a parameter-free formula and
  $M\sat \forall a\,(F(a)\leftrightarrow a=\nu(x))$ (length of $F$ below
  the bound).
\end{description}
So both \emph{satisfaction} ($M\sat F$) and \emph{definability}
($F$ defines a value in $M$) are formulas of the language, evaluated at a
model-object~$M$.  Nothing here is meta.

\subsection{The uniformity is models-inside-models}

This is the step the summary hides.  Krivine does \emph{not} argue
``$\nu(n)^M=\nu(n)^U$ for all external models $M,U$'' as an abstract
invariance.  He argues internally:

\begin{enumerate}
\item Take a model $M\sat\Thy$.  By the contradiction hypothesis
  $M\sat\Con(\Thy)$.
\item Since $M\sat\Con(\Thy)$, \emph{inside $M$} there is an object $N'$
  which $M$ believes is a model of $\Thyo$.  (Krivine cites his own
  earlier ``consistent $\Rightarrow$ has a model,'' p.~123/125,
  \emph{applied internally to $M$}.)
\item The two models are glued by the internal soundness/completeness
  equivalence:
  \[
     \text{``}N'\sat F\text{'' is true in } M
        \quad\Longleftrightarrow\quad F\in\Thy_M .
  \]
  (For every closed $F$ with parameters in $M$.)  This is what forces
  agreement of the Berry values.
\item Hence $\nu(n)^{M}\ \ge\ \nu(n)^{N'}$, and the symmetric argument
  gives the reverse, so $\nu(n)$ is \emph{model-independent} and bounded:
  $\nu(n)^{M}\le 2^{2^{2m}}$ (an explicit exponential tower from the
  counting; the exact height is not fully legible).
\end{enumerate}

\begin{remark}[where $\Con$ is consumed]\label{rmk:con}
The hypothesis $\Thy\vdash\Con(\Thy)$ is used at \emph{exactly one place}:
step~(2), the passage $M\rightsquigarrow N'$.  Consistency is what lets
$M$ produce an internal model $N'$ to compare against.  Everything else
(counting, length, the self-reference) is hypothesis-free.
\end{remark}

\subsection{The counting bound}

The number of parameter-free formulas of length~$<l$ is controlled by a
primitive-recursive recurrence.  As legibly as the scan permits:
\[
   p_0 = 2^{2m},\qquad
   p_l \ \le\ p_{l-1} + 1 + 2^{m}\,p_{l-1},
\]
where $p_l$ counts parameter-free formulas of length~$<l$ over the
$2^{2m}$ first symbols.  This yields the explicit bound
$\nu(m_0)^M\le 2^{2^{2m}}$ used above.  (Constants reported as printed;
the BRA translation only needs ``$p_l$ is a closed-form
primitive-recursive function with an exponential majorant.'')

\subsection{The Berry self-reference}

Finally (p.~127, ``on montrera un peu plus loin''):
\begin{quote}
one can choose an \emph{intuitive} (meta) integer $m_0$ and write a
parameter-free statement $D(x)$ equivalent to ``$x=\nu(m_0)$'' that is
itself of length~$<m_0$.
\end{quote}
Then $\nu(m_0)$ \emph{is} definable by a formula of length~$<m_0$,
contradicting its definition as the least integer not so definable.
This refutes $\Thy\vdash\Con(\Thy)$.

\begin{remark}[why this is not Berry's paradox]
The self-reference is harmless as a paradox because $\nu$ is defined
through \emph{provability/satisfaction in a model} ($M\sat\dots$), not
through raw evaluation: deciding ``$x=\nu(m_0)$'' asks whether the model
satisfies a bounded statement, not for $\nu$ to run itself.  This is the
same reason the $\Sigma_1$ (provable-naming) version is sound where the
$\Sigma_0$ (Eval-naming) version of \texttt{berry-krivine-design.tex}
was not.
\end{remark}

% =================================================================
\section{The engine, abstracted}
% =================================================================

Stripped of $ZF$, Krivine's proof has four moving parts:
\begin{description}
\item[(M) a canonical model] in which satisfaction $M\sat F$ is a
  predicate of the language;
\item[(I) an internal model $N'$] obtained \emph{from $\Con$} inside
  $M$, with the soundness equivalence $N'\sat F \Leftrightarrow F\in\Thy$;
\item[(C) a counting/length calculus] giving an explicit exponential
  bound on $\nu$;
\item[(B) a length self-reference] $D(x)\equiv$``$x=\nu(m_0)$'' of
  length~$<m_0$, firing the Berry contradiction.
\end{description}
Parts (M)+(I) are von~Neumann's ``copy the antinomy directly'' route:
the contradiction is built from the model and its internal model, not
from a G\"odel fixed point.  Parts (C)+(B) are the Berry antinomy proper.
Note that \emph{full completeness is never invoked}; the only direction
used is ``$\Con\Rightarrow$ internal model'' (the easy half, internally),
plus soundness $N'\sat F\Rightarrow F\in\Thy$.

% =================================================================
\section{Translation to Guard's BRA}
% =================================================================

We use the BRA notion of model settled in earlier discussion: a model is
an enumerator $M:\Fun_1$ with
\[
   M\sat a \ :=\ \exists z.\;M(z)=a,
\]
and the \emph{canonical} model is the verifier $\thmT$ (theorem
enumerator).  Internal models are $\thmT$-relativisations built with the
substitution functor $\sub$.  The dictionary:

\begin{center}
\begin{tabular}{p{0.46\textwidth}p{0.46\textwidth}}
\hline
\textbf{Krivine (pp.~126--127)} & \textbf{BRA-internal} \\
\hline
outer model $M\sat\Thy$
  & $\thmT$; \quad $M\sat a := \exists z.\,\thmT(z)=a$
    (range-membership, $\Sigma_1$, Skolemised as a function --- the
    \texttt{big\_term} pattern of the shipped proof) \\[2pt]
inner model $N'$ inside $M$, from $\Con$
  & \emph{not needed as a literal object.}  Its only content is nested
    provability: $\thmT$ reasoning about $\thmT$.  A $\sub$-relativisation
    of $\thmT$ would mirror Krivine, but is optional faithfulness, not a
    requirement (see {\S}\ref{sec:why-no-reloc}) \\[2pt]
soundness link ``$N'\sat F$ in $M$ iff $F\in\Thy$''
  & $\thmT$'s validating-decoder invariant $+$ $\Sigma_0$-completeness
    (closed-term evaluation, the \texttt{numEq}/eval chain) \\[2pt]
uniformity $\nu(n)^M=\nu(n)^{N'}$
  & \emph{provable $\Sigma_1$-completeness} $\Box\varphi\to\Box\Box\varphi$
    (HBL\,D3): $\thmT$ proves its own definability judgements.  This is
    the syntactic residue of ``inner model agrees with outer'' \\[2pt]
counting recurrence $p_l$
  & internal length/counting lemmas (compressed numerals so $\Len$ is
    $O(\log)$) \\[2pt]
meta $m_0$, $D(x)\equiv x{=}\nu(m_0)$, $\Len<m_0$
  & the length-accounting self-reference (the residual inequality) \\[2pt]
hypothesis $\Thy\vdash\Con(\Thy)$
  & \texttt{Deriv ConSchema}, consumed exactly at $\thmT\!\rightsquigarrow\!$
    inner model (Remark~\ref{rmk:con}) \\
\hline
\end{tabular}
\end{center}

\paragraph{The key correction over the companion note.}  What
\texttt{berry-krivine-design.tex} {\S}``Status'' called the open
``definability cut / Lindenbaum cut'' is, in Krivine's \emph{semantic}
proof, the internal model $N'$.  But $N'$ is not load-bearing in BRA: it
is the \emph{semantic face of nested provability}.  Replacing
``definable in a model'' by ``$\thmT$-provable,'' the cut dissolves into
a single derivability lemma about $\thmT$ (provable $\Sigma_1$-
completeness, {\S}\ref{sec:why-no-reloc}).  So the old ``open cut'' is
neither open nor a cut: it is $\thmT$'s D3, repackaged.

\paragraph{$ZF$ vs.\ BRA.}  Krivine works over $ZF$, where a model is a
\emph{set-object} $N'$ and $\forall F$, $\exists M$ range over the
universe.  BRA is quantifier-free primitive-recursive arithmetic, so:
\begin{itemize}
\item ``definable in a model'' becomes ``$\thmT$-provable'' (the only
  expressible surrogate; see {\S}\ref{sec:why-no-reloc});
\item the quantifiers $\forall F$ / $\exists M$ become \emph{function
  building} (Skolemisation), exactly as Guard's $\Sigma_1$ provability is
  realised by constructing \texttt{big\_term} rather than by an internal
  $\exists$;
\item ``$M\sat F$'' becomes range-membership of the enumerator $\thmT$.
\end{itemize}
The counting/length calculus is unchanged --- it is finitary and
primitive-recursive in either setting.

% =================================================================
\section{What the model really provides: bivalence, not decidability}
\label{sec:why-no-reloc}
% =================================================================

\emph{Correction of an earlier draft.}  We first claimed Krivine's
definability is \emph{decidable} (satisfaction in a set-model), and that
BRA --- lacking set-models --- therefore could not host the argument.
That is wrong.  A model of arithmetic is \emph{infinite}, so
``$M\sat\forall w\,(F(w)\leftrightarrow w=p)$'' is a universal over an
infinite domain, not a finite check.  Krivine's definability is
\textbf{not decidable} either.

What a model actually supplies is \textbf{bivalence}: inside a model $U$,
every statement --- including the undecidable definability ones --- has a
definite truth value, so ``$\nu=$ least $p$ with [property]'' is
well-defined \emph{regardless of decidability}.  $\nu$ is pinned by
\emph{working inside a model}, not by being computed.

The syntactic counterpart of bivalence is \textbf{provability plus the
consistency hypothesis}.  Under the translation ``true in $U$''
$\longrightarrow$ ``provable in $\thmT$'' ($\Sigma_1$), the bivalent value
of $\nu$ becomes: $\nu$ \emph{handled at the formula level} via the Berry
formula $B$ together with provability and $\Con$ --- \emph{never
computed}.  Definability stays undecidable throughout, exactly as in
Krivine; the models-inside-models / uniformity step becomes nested
provability $\Box\varphi\to\Box\Box\varphi$ (provable
$\Sigma_1$-completeness, HBL\,D3).

\textbf{Conclusion: Krivine's argument does render syntactically.}  No
decidability is needed --- it was never available, not even in Krivine ---
and no relativised model is needed; what is needed is
$\thmT + \Con + {}$D3.  (This supersedes two earlier detours: an
over-optimistic ``build a relativised model for decidable definability''
and an over-pessimistic ``Krivine cannot adapt.''  The stable statement:
it renders via $\Sigma_1$-provability, definability is undecidable, and
$\nu$ lives at the formula level.)

\paragraph{Complexity, exactly.}  ``$p$ is $n$-definable'' is $\Sigma_1$
(a bounded $\exists F$ over the $\Sigma_1$ predicate ``$\thmT$ proves $F$
defines $p$'').  ``$p$ is the \emph{least} non-$n$-definable'' is
$\neg\mathrm{Def}_n(p)\,(\Pi_1)\,\wedge\,\forall q<p\,\mathrm{Def}_n(q)\,
(\Sigma_1)$ --- a $\Pi_1\wedge\Sigma_1$ conjunction, \emph{not} $\Sigma_1$.
It cannot be $\Sigma_1$: $\nu$ is total (by the counting) but not
computable, and a total function with a $\Sigma_1$ graph would be
computable.  The $\Pi_1$ conjunct is exactly where $\Con$ is spent; the
$\Sigma_1$ conjunct is the witness-producing part.

% =================================================================
\section{What is left of the argument: a model-free Berry proof}
\label{sec:whatleft}
% =================================================================

With the relativisation gone, the argument is precisely \textbf{Berry's
paradox with $\thmT$ as the definability oracle} --- no models, no cut.
The skeleton:

\begin{enumerate}
\item \textbf{Definability via $\thmT$ ($\Sigma_1$).}  ``$p$ is
  $n$-definable'' $:=$ $\exists F$ ($\Len(F)<n$ and $\thmT$ proves
  $\forall a\,(F(a)\leftrightarrow a=\mathbf{k}_p)$).  Skolemised as a
  function (the \texttt{big\_term} pattern).
\item \textbf{Counting.}  $<2^{2n}$ formulas of length $<n$, hence
  $<2^{2n}$ $n$-definable numbers; so some number $\le 2^{2n}$ is not
  $n$-definable.
\item \textbf{Berry function.}  $\nu(n) := $ least non-$n$-definable
  number ($\le 2^{2n}$).  It is total but \emph{not computable} (graph
  $\Pi_1\wedge\Sigma_1$, {\S}\ref{sec:why-no-reloc}), so it is \emph{not}
  a $\Fun_1$ and is never computed --- it is carried by the formula $B$
  and reasoned about via provability $+\,\Con$.
\item \textbf{Self-reference by length.}  Let $B(x,y)\equiv$ ``$x$ is the
  least non-$y$-definable number''; $\thmT$ is a \emph{fixed} $\Fun_1$,
  so $B$'s text is constant-size and $\Len(B(x,\mathbf{k}_{m_0})) =
  c + |\mathbf{k}_{m_0}|$.  With compressed numerals $|\mathbf{k}_{m_0}| =
  O(\log m_0)$, so for a concrete meta $m_0$, $\Len(B(\cdot,m_0))<m_0$.
\item \textbf{Contradiction (where $\Con$ and D3 are spent).}
  \begin{itemize}
  \item $\Con(\Thy)$ $\Rightarrow$ $\thmT$ proves the \emph{negative}
    fact ``$\nu(m_0)$ is \emph{not} $m_0$-definable'' (proving a
    non-provability needs consistency);
  \item provable $\Sigma_1$-completeness (D3) $\Rightarrow$ $\thmT$ proves
    the \emph{positive} facts ``each $p<\nu(m_0)$ \emph{is}
    $m_0$-definable.''
  \end{itemize}
  Together $\thmT$ proves ``$B(\cdot,m_0)$ defines $\nu(m_0)$'' with
  $\Len(B(\cdot,m_0))<m_0$ --- i.e.\ $\nu(m_0)$ \emph{is}
  $m_0$-definable, contradicting step~3.  Hence $\Thy\not\vdash\Con(\Thy)$.
\end{enumerate}

\paragraph{Why this is the \emph{fix}, not the rejected version.}  The
companion note rejected a Berry attempt --- but a \emph{different} one.
The rejected version used \textbf{Eval-definability} ($\Sigma_0$): ``$p$
is named by $F$'' $:=$ \emph{running} $F$ yields $p$.  To decide that,
the deciding function must run $\Eval$ on every short code including its
own (since $\Len(\code\beta)<n_0$), an infinite self-simulation; $\beta$
is not total, ``$\beta$ defines $\nu$'' is \emph{false} in $\mathbb N$,
and the contradiction never fires.  That is the classical dissolution of
Berry's paradox.

The version above uses \textbf{$\thmT$-definability} ($\Sigma_1$):
``$p$ is named by $F$'' $:=$ $\thmT$ \emph{proves} $F$ defines $p$.  This
asks for the \emph{existence of a derivation} (a search), never for
$\beta$ to run itself; and $\thmT$ appears in $B$ only \emph{by name}, as
a constant-size sub-term, never unfolded.  So there is no regress, $\nu$
is a genuine bounded number, and the self-reference produces a real
theorem.  The two differ exactly as $\Sigma_0$ (compute) differs from
$\Sigma_1$ (prove) --- which is the entire content of the companion
note's ``Status'' correction.  \emph{Berry-with-$\thmT$ is the endorsed
route; Berry-with-$\Eval$ is the dead one.}

% =================================================================
\section{Expressing it in Guard's quantifier-free setting}
\label{sec:guard}
% =================================================================

Guard's BRA is \emph{quantifier-free}: object formulas are atomic
equations, negation, implication, and \emph{a free variable is an
implicit universal}.  The two complexity levels above are therefore
carried by two different mechanisms --- the \emph{same} two the shipped
diagonal proof already uses.

\paragraph{$\Pi_1$ / universal $\to$ open formulas.}  ``$x$ is not
short-definable'' $=\neg\mathrm{DefWit}(F,\pi,x,m_0)$ with $F,\pi$ free
is, in Guard, already $\forall F,\pi\,(\dots)$.  This is ``the language
has only universal statements.''  $\Con$ itself is this form:
\texttt{neg (eqF (ap1 thmT (var 0)) codeFalse)}, i.e.\
$\forall v.\,\thmT(v)\neq\mathrm{codeFalse}$.

\paragraph{$\Sigma_1$ $\to$ Skolemised witness terms, not formulas.}
``$x$ is short-definable'' is $\Sigma_1$ and \emph{cannot} be an object
formula (one cannot recover $\exists$ by negating an open $\forall$).  It
is realised exactly as the shipped proof realises provability: by
\emph{building} the witness $(F,\pi)$ as a term and asserting the
\emph{decidable} equation $\mathrm{DefWit}$ --- the \texttt{big\_term}
technique.

\paragraph{The contradiction, in Guard's shape.}  Signature
\texttt{Deriv ConSchema -> Deriv falseF}, as in \texttt{Thm14GodelII}.
Build a short $B$; under the $\Con$ hypothesis, build the witness that
$\thmT$ proves ``$B$ defines its value $v$'' --- this build consumes
$\Con$ (for the $\Pi_1$ ``$v$ not short-definable'') and uses an
\emph{induction-built family} of witnesses (for the $\Sigma_1$ ``every
$q<v$ is short-definable'').  Then $B$ short $\wedge$ ``$B$ defines $v$''
gives $v$ short-definable, while $\thmT\vdash B(v)$ gives $v$ \emph{not}
short-definable; so $\thmT\vdash\bot$, i.e.\ \texttt{Deriv falseF}.

\paragraph{The Guard-specific friction.}  $\nu$ is non-computable and $B$
is irreducibly $\Pi_1\wedge\Sigma_1$, so --- unlike in PA --- $B$ is
\emph{not} a single quantifier-free object formula.  Its $\Pi_1$ part is
an open formula; its $\Sigma_1$ ``least / all-smaller-definable'' part
must be carried as a \emph{bound-parameterised family of Skolemised
witnesses, built by} \texttt{ruleIndNat}.  That induction-built family is
the ``D3 over a range'' component.  It is heavier than PA's ``just write
the formula,'' because quantifier-freeness forces the $\Sigma_1$ content
out of the object language into Deriv-level term-building --- but it is
precisely what \texttt{big\_term} $+$ \texttt{ruleIndNat} were built for.

\paragraph{The full Guard recipe.}
\begin{enumerate}
\item \texttt{Len} $+$ compressed numerals (so
  $\Len(B)=c+O(\log m_0)<m_0$) --- new;
\item $\mathrm{DefWit}$ as a decidable equation (verify a proof, check
  length) --- routine;
\item the counting bound \texttt{\#definable <= \#short-formulas} with
  \textbf{$\Con$-injectivity} (a formula names $\le 1$ number unless
  inconsistent) --- the hard internal combinatorics;
\item open-formula $\Pi_1$ non-definability $+$ \texttt{ruleInst} on
  $\Con$ --- reuse the shipped idiom;
\item the \texttt{ruleIndNat}-built witness family for ``all smaller
  definable'' --- the mixed-$B$ workaround;
\item assemble \texttt{Deriv ConSchema -> Deriv falseF}.
\end{enumerate}
The genuine unknowns are (3) and (1); neither is a wall.

% =================================================================
\section{Honest accounting}
% =================================================================

\paragraph{Routine (methods already in the codebase).}
\begin{itemize}
\item $\Pi_1$/universal statements as open formulas; $\Sigma_1$
  definability as Skolemised witness terms (the \texttt{big\_term}
  pattern of the shipped \texttt{Thm14GodelII}) --- the two idioms of
  {\S}\ref{sec:guard};
\item $\thmT$ as the single provability oracle (no relativised copy,
  {\S}\ref{sec:why-no-reloc});
\item $\Sigma_0$-completeness for the decidable $\mathrm{DefWit}$ check
  ($=$ verifying a given proof, $=$ numeral evaluation);
\item $\Con$ as hypothesis (\texttt{Deriv ConSchema}), consumed via
  \texttt{ruleInst} at the $\Pi_1$ non-definability step.
\end{itemize}

\paragraph{Genuinely remaining (no new concept, real engineering).}
\begin{enumerate}
\item the internal counting recurrence $p_l$ and its exponential
  majorant, with compressed numerals so the length bookkeeping closes
  (the recommended ``gate experiment'');
\item the concrete meta integer $m_0$ and the inequality
  $\Len(B(\cdot,m_0))<m_0$ for the Berry formula (step~4 of
  {\S}\ref{sec:whatleft}) --- the one substantive step, now
  \emph{well-posed} ($B$ is the $\thmT$-definability predicate at the
  fixed bound $m_0$, not a vague ``nameability'');
\item provable $\Sigma_1$-completeness (D3) for $\thmT$, for the positive
  definability facts in step~5 --- the same derivability ingredient the
  diagonal/L\"ob proof uses, repackaged.
\end{enumerate}

\paragraph{Relation to the shipped proof.}  This is a genuinely
different proof of the same theorem, not a re-skin: the shipped
\texttt{Thm14GodelII} uses the G\"odel sub-diagonal fixed point, whereas
this route uses Berry counting $+$ $\Sigma_1$-provability definability ---
von~Neumann's direct, antinomy-copying method, with no fixed point.  Both
consume \texttt{Deriv ConSchema} and produce \texttt{Deriv falseF}, by
different mechanisms.

\paragraph{Status of the companion note.}  This note supersedes the
{\S}``Status'' diagnosis of \texttt{berry-krivine-design.tex}: the open
step is neither an unspecified ``Lindenbaum cut'' nor a $\sub$-relativised
inner model.  It is the model-free $\Sigma_1$ Berry recipe of
{\S}\ref{sec:guard}, with $\Con$ consumed at the $\Pi_1$ non-definability
fact.  The two genuine pieces of work are the $\Con$-injective counting
and the length self-reference; D3 supplies the ``all-smaller-definable''
family.

% =================================================================
\section*{Sources}
% =================================================================
\begin{itemize}
\item J.-L. Krivine, \emph{Th\'eorie axiomatique des ensembles},
  pp.~126--127 (photographs \texttt{IMG\_0196}, \texttt{IMG\_0197}).
\item Companion design note: \texttt{berry-krivine-design.tex} (Berry/
  naming components, $\Sigma_0$-vs-$\Sigma_1$ analysis).
\item Shipped formalisation: \texttt{BRA4/Thm/Thm14GodelII.agda}
  (diagonal route), and the model notion ($\thmT$, $\sub$,
  $\Sigma_0$-completeness) developed there.
\item C.~Guard, \emph{Automated logic and the foundations of recursive
  arithmetic} (1963); A.~Church, \emph{Introduction to Mathematical
  Logic} (1957).
\end{itemize}

\end{document}
